\pdfminorversion=7
\documentclass[spanish,aspectratio=1610,xcolor={dvipsnames,table}]{beamer}

% Lienzo amplio conservando proporción 16:10
% Original: 19.2cm x 12cm
\geometry{paperwidth=28cm,paperheight=17.5cm}

\usepackage[utf8]{inputenc}
\usepackage[T1]{fontenc}
\usepackage[spanish,es-nodecimaldot]{babel}
\usepackage{lmodern}
\usepackage{amsmath}
\usepackage{tikz}
\usetikzlibrary{
  arrows.meta,
  babel,
  calc,
  angles,
  quotes,
  positioning,
  decorations.pathreplacing
}

% =============================================================
% Plantilla genérica: sin datos personales, institucionales ni secciones.
% Título del documento: Diagramas TikZ.
% Ejemplo central: plano inclinado con aparición paso a paso.
% =============================================================

\definecolor{navy}{HTML}{163B5C}
\definecolor{blue}{HTML}{175CD3}
\definecolor{red}{HTML}{B42318}
\definecolor{green}{HTML}{067647}
\definecolor{orange}{HTML}{B54708}
\definecolor{paper}{HTML}{F6F8F7}
\definecolor{ink}{HTML}{1E2A32}
\definecolor{muted}{HTML}{667085}

\mode<presentation>{
  \usetheme{default}
  \usefonttheme{professionalfonts}
  \setbeamertemplate{navigation symbols}{}
  \setbeamercovered{invisible}
}

\setbeamercolor{background canvas}{bg=white}
\setbeamercolor{normal text}{fg=ink,bg=white}
\setbeamercolor{structure}{fg=blue}
\setbeamercolor{frametitle}{fg=white,bg=navy}
\setbeamerfont{frametitle}{size=\large,series=\bfseries}

\setbeamertemplate{frametitle}{%
  \nointerlineskip
  \begin{beamercolorbox}[wd=\textwidth,ht=1.00cm,dp=0.22cm,leftskip=0.42cm,rightskip=0.25cm]{frametitle}
    \insertframetitle
  \end{beamercolorbox}
  {\color{orange}\rule{\textwidth}{1.2pt}}
}

\setbeamertemplate{footline}{%
  \leavevmode%
  \hbox{%
    \begin{beamercolorbox}[wd=.80\paperwidth,ht=2.5ex,dp=1ex,leftskip=1em]{author in head/foot}%
      \color{white}Diagramas TikZ
    \end{beamercolorbox}%
    \begin{beamercolorbox}[wd=.20\paperwidth,ht=2.5ex,dp=1ex,rightskip=1em plus 1fil]{date in head/foot}%
      \color{white}\hfill\insertframenumber/\inserttotalframenumber
    \end{beamercolorbox}%
  }%
}
\setbeamercolor{author in head/foot}{bg=navy}
\setbeamercolor{date in head/foot}{bg=navy}

\title{Diagramas TikZ}
\date{}

\begin{document}

% -------------------------------------------------------------
% Portada con formato conservado
% -------------------------------------------------------------
\begin{frame}[plain]
	\begin{tikzpicture}[remember picture,overlay]
		\fill[navy] (current page.south west) rectangle (current page.north east);
		\fill[blue,opacity=0.35] ([xshift=0.58\paperwidth]current page.south west) rectangle (current page.north east);

		\fill[orange]
		([xshift=1.15cm,yshift=2.0cm]current page.south west)
		rectangle
			([xshift=1.35cm,yshift=7.3cm]current page.south west);

		\fill[green]
		([xshift=1.55cm,yshift=2.0cm]current page.south west)
		rectangle
			([xshift=1.75cm,yshift=6.4cm]current page.south west);

		\node[
			anchor=west,
			text=white,
			font=\Huge\bfseries
		] at ([xshift=2.25cm,yshift=1.0cm]current page.west) {Diagramas TikZ};

		\node[
			anchor=west,
			text=white!75,
			font=\large
		] at ([xshift=2.28cm,yshift=0.25cm]current page.west) {Ejemplo con plano inclinado y overlays};
	\end{tikzpicture}
\end{frame}

% ------------------------------------------------------------------
% Diapositiva principal
% ------------------------------------------------------------------

\begin{frame}{Carrito en plano inclinado}
	\begin{center}
		\begin{tikzpicture}[
				caja/.style={draw, rounded corners, fill=blue!10, inner sep=6pt},
				flecha/.style={-{Latex}, thick},
				nota/.style={draw, rounded corners, fill=yellow!20, align=center}
			]

			\node[caja] (eq1) {$2x + 3 = 11$};
			\node[caja, below=1.2cm of eq1] (eq2) {$2x = 11 - 3$};
			\node[caja, below=1.2cm of eq2] (eq3) {$2x = 8$};
			\node[caja, below=1.2cm of eq3] (eq4) {$x = \frac{8}{2}$};
			\node[caja, below=1.2cm of eq4] (eq5) {$x = 4$};

			\draw[flecha] (eq1) -- node[right] {Restar 3 en ambos lados} (eq2);
			\draw[flecha] (eq2) -- node[right] {Simplificar} (eq3);
			\draw[flecha] (eq3) -- node[right] {Dividir entre 2} (eq4);
			\draw[flecha] (eq4) -- node[right] {Resultado} (eq5);

			\node[nota, right=2.5cm of eq3] (nota1) {El objetivo es\\dejar sola a $x$};
			\draw[flecha, dashed] (nota1) -- (eq3);

		\end{tikzpicture}
	\end{center}
\end{frame}

% -------------------------------------------------------------------
% Diapositiva: plano inclinado con DCL, modelo dinámico y cinemática
% -------------------------------------------------------------------
\begin{frame}{Plano inclinado con DCL, modelo dinámico y cinemática}

	\centering

	\begin{tikzpicture}[
			x=1cm,
			y=1cm,
			>=Stealth,
			force/.style={-{Stealth[length=3.2mm,width=2.4mm]}, very thick, line cap=round},
			component/.style={-{Stealth[length=2.8mm,width=2.1mm]}, very thick, dashed, line cap=round},
			axis/.style={-{Stealth[length=2.6mm,width=2.0mm]}, thick, muted},
			panel/.style={rounded corners=4pt, draw=muted!35, fill=paper, inner sep=5pt},
			label/.style={font=\scriptsize, inner sep=1pt},
			smalllabel/.style={font=\tiny, inner sep=1pt},
			eq/.style={font=\footnotesize, align=left, text=ink},
			every node/.style={font=\small}
		]

		% =============================================================
		% DISTRIBUCIÓN GENERAL
		% =============================================================
		% Panel izquierdo: Sistema físico
		% Panel central: DCL
		% Panel derecho superior: Modelo dinámico
		% Panel derecho inferior: Cinemática
		% =============================================================

		% -------------------------------------------------------------
		% Panel 1: sistema físico más alto
		% -------------------------------------------------------------
		\only<1->{
				\node[
					panel,
					minimum width=12.2cm,
					minimum height=9.15cm,
					anchor=north west
				] at (-13.2,6.85) {};

				\node[
					anchor=west,
					font=\bfseries\small,
					text=navy
				] at (-12.75,6.38) {Sistema físico};
			}

		% Coordenadas del plano
		\def\ang{28}
		\coordinate (A) at (-12.35,-0.55);
		\coordinate (B) at (-2.15,-0.55);
		\coordinate (C) at (-2.15,4.55);
		\coordinate (P) at ($(A)!0.56!(C)$);
		\coordinate (Pc) at ($(P)+(0.00,0.35)$);

		% Plano inclinado
		\only<1->{
				\fill[gray!11] (A) -- (B) -- (C) -- cycle;
				\draw[thick, muted] (A) -- (B) -- (C);
				\draw[very thick, ink] (A) -- (C);

				% Ángulo theta
				\draw[thick, ink!70]
				($(A)+(0.92,0)$)
				arc[start angle=0,end angle=\ang,radius=0.92];

				\node[label] at ($(A)+(1.12,0.30)$) {$\theta$};

				% Etiqueta horizontal
				\node[label,text=muted] at ($(A)!0.50!(B)+(0,-0.34)$) {horizontal};

				% Etiqueta del plano, bajada y movida para no encimarse con W sin(theta)
				\node[
					label,
					text=muted,
					rotate=\ang
				] at ($(A)!0.22!(C)+(0.25,-0.52)$) {plano inclinado};
			}

		% Bloque
		\only<2->{
				\begin{scope}[shift={(P)}, rotate=\ang]
					\draw[
						fill=blue!12,
						draw=blue!80!black,
						thick,
						rounded corners=1pt
					] (-0.62,-0.34) rectangle (0.62,0.42);

					\node[text=blue!45!black] at (0,0.04) {$m$};
				\end{scope}
			}

		% Ejes locales
		\only<2->{
				\coordinate (E) at ($(P)+(1.55,0.90)$);

				\draw[axis] (E) -- ++({1.35*cos(\ang)},{1.35*sin(\ang)})
				node[right,label,text=ink] {$+x$};

				\draw[axis] (E) -- ++({1.08*cos(\ang+90)},{1.08*sin(\ang+90)})
				node[above,label,text=ink] {$+y$};
			}

		% Peso W
		\only<3->{
				\draw[force, red] (Pc) -- ++(0,-2.35)
				node[below=2pt,label,text=red] {$W=mg$};
			}

		% Normal
		\only<4->{
				\draw[force, blue] (Pc) -- ++({1.75*cos(\ang+90)},{1.75*sin(\ang+90)})
				node[above left=1pt,label,text=blue] {$N$};
			}

		% Fuerza paralela al plano
		\only<5->{
				\draw[force, green] (Pc) -- ++({2.45*cos(\ang)},{2.45*sin(\ang)})
				node[right=2pt,label,text=green] {$F$};
			}

		% Componentes del peso
		\only<6->{
				% Se desplazan ligeramente para no encimarse con W
				\coordinate (Pcomp) at ($(Pc)+(-0.18,-0.08)$);

				\draw[component, orange] (Pcomp) -- ++({-1.95*cos(\ang)},{-1.95*sin(\ang)})
				node[left=2pt,label,text=orange] {$W\sin\theta$};

				\draw[component, orange] (Pcomp) -- ++({-1.55*cos(\ang+90)},{-1.55*sin(\ang+90)})
				node[below left=1pt,label,text=orange] {$W\cos\theta$};

				% Nota de descomposición, arriba a la izquierda y sin tocar las flechas
				\node[
					eq,
					anchor=north west,
					text width=5.0cm
				] at (-12.55,5.95) {
					\textbf{Descomposición del peso:}\\[-1mm]
					\[
						\begin{aligned}
							W_x & = W\sin\theta \\
							W_y & = W\cos\theta
						\end{aligned}
					\]
				};
			}

		% -------------------------------------------------------------
		% Panel 2: DCL más alto
		% -------------------------------------------------------------
		\only<7->{
				\node[
					panel,
					minimum width=6.2cm,
					minimum height=9.15cm,
					anchor=north west
				] at (-0.45,6.85) {};

				\node[
					anchor=west,
					font=\bfseries\small,
					text=navy
				] at (-0.05,6.38) {DCL: cuerpo libre};

				\coordinate (D) at (2.65,2.00);

				% Bloque DCL
				\draw[
					fill=blue!12,
					draw=blue!80!black,
					thick,
					rounded corners=1pt
				] ($(D)+(-0.46,-0.34)$) rectangle ($(D)+(0.46,0.34)$);

				\node[text=blue!45!black] at (D) {$m$};

				% Ejes del DCL
				\draw[axis] ($(D)+(-2.05,-2.35)$) -- ++(1.55,0)
				node[right,label,text=ink] {$+x$};

				\draw[axis] ($(D)+(-2.05,-2.35)$) -- ++(0,1.45)
				node[above,label,text=ink] {$+y$};

				% Fuerzas DCL
				\draw[force, green] (D) -- ++(2.00,0)
				node[right=2pt,label,text=green] {$F$};

				\draw[force, blue] (D) -- ++(0,2.00)
				node[above=2pt,label,text=blue] {$N$};

				\draw[force, orange] (D) -- ++(-2.00,0)
				node[left=2pt,label,text=orange] {$W\sin\theta$};

				\draw[force, orange] (D) -- ++(0,-2.00)
				node[below=2pt,label,text=orange] {$W\cos\theta$};
			}

		% -------------------------------------------------------------
		% Panel 3: modelo dinámico más alto
		% -------------------------------------------------------------
		\only<8->{
				\node[
					panel,
					minimum width=8.6cm,
					minimum height=6.15cm,
					anchor=north west
				] at (6.25,6.85) {};

				\node[
					anchor=west,
					font=\bfseries\small,
					text=navy
				] at (6.65,6.38) {Modelo dinámico};

				\node[
					eq,
					anchor=north west,
					text width=7.70cm
				] at (6.70,5.78) {
					\[
						\begin{aligned}
							\sum F_x & = F-W\sin\theta   \\[2mm]
							\sum F_y & = N-W\cos\theta=0
						\end{aligned}
					\]
				};
			}

		\only<9->{
				\node[
					eq,
					anchor=north west,
					text width=7.70cm
				] at (6.70,3.70) {
					Como $W=mg$:
					\[
						ma=F-mg\sin\theta
					\]

					\[
						\boxed{
							a=\dfrac{F-mg\sin\theta}{m}
						}
					\]
				};
			}

		% -------------------------------------------------------------
		% Panel 4: cinemática más abajo y con aire
		% -------------------------------------------------------------
		\only<10->{
				\node[
					panel,
					minimum width=8.6cm,
					minimum height=3.00cm,
					anchor=north west
				] at (6.25,0.25) {};

				\node[
					anchor=west,
					font=\bfseries\small,
					text=navy
				] at (6.65,-0.22) {Cinemática};

				\node[
					eq,
					anchor=north west,
					text width=7.70cm
				] at (6.70,-0.72) {
					Si $v_0=0$ y $t=3\,s$:
					\[
						\boxed{x=\frac{1}{2}a(3)^2}
						\qquad
						\boxed{v_f=a(3)}
					\]
				};
			}

	\end{tikzpicture}
\end{frame}

% -------------------------------------------------------------
% Frame 1: construcción de vectores y componentes
% -------------------------------------------------------------
\begin{frame}{Sumatoria de vectores: construcción y componentes}

	\centering

	\begin{tikzpicture}[
			x=1cm,
			y=1cm,
			>=Stealth,
			vector/.style={-{Stealth[length=3.2mm,width=2.4mm]}, very thick, line cap=round},
			component/.style={-{Stealth[length=2.8mm,width=2.1mm]}, very thick, dashed, line cap=round},
			axis/.style={-{Stealth[length=2.6mm,width=2.0mm]}, thick, muted},
			panel/.style={rounded corners=4pt, draw=muted!35, fill=paper, inner sep=5pt},
			label/.style={font=\scriptsize, inner sep=1pt},
			eq/.style={font=\footnotesize, align=left, text=ink},
			every node/.style={font=\small}
		]

		% =============================================================
		% Datos del problema
		% F1 = 50, theta1 = 60°
		% F2 = 70, theta2 = 113°
		%
		% Fx1 = 25.00
		% Fy1 = 43.30
		% Fx2 = -27.35
		% Fy2 = 64.44
		% =============================================================

		% -------------------------------------------------------------
		% Panel izquierdo: sistema vectorial
		% -------------------------------------------------------------
		\only<1->{
				\node[
					panel,
					minimum width=15.0cm,
					minimum height=9.25cm,
					anchor=north west
				] at (-13.15,6.85) {};

				\node[
					anchor=west,
					font=\bfseries\small,
					text=navy
				] at (-12.70,6.35) {Sistema vectorial};
			}

		% Escala y puntos
		\coordinate (O) at (-7.00,0.35);
		\def\s{0.052}

		\coordinate (Fone) at ($(O)+({50*cos(60)*\s},{50*sin(60)*\s})$);
			\coordinate (Ftwo) at ($(O)+({70*cos(113)*\s},{70*sin(113)*\s})$);

			\coordinate (FxOne) at ($(O)+({25.00*\s},0)$);
			\coordinate (FxTwo) at ($(O)+({-27.35*\s},0)$);

			% Ejes
			\only<1->{
			\draw[axis] ($(O)+(-4.6,0)$) -- ($(O)+(5.2,0)$)
			node[right,label,text=ink] {$+x$};

			\draw[axis] (O) -- ++(0,5.2)
			node[above,label,text=ink] {$+y$};

			\node[label,text=muted] at ($(O)+(-0.20,-0.25)$) {$O$};
			}

			% Vector F1
			\only<2->{
			\draw[vector, blue] (O) -- (Fone)
			node[pos=0.64, below right=2pt, label, text=blue] {$\vec{F}_1$};

			\node[
				label,
				text=blue,
				anchor=west
			] at ($(Fone)+(0.18,0.06)$) {$F_1=50$};
			}

			% Ángulo theta1
			\only<3->{
			\draw[blue, thick]
			($(O)+(0.95,0)$)
			arc[start angle=0,end angle=60,radius=0.95];

			\node[label,text=blue] at ($(O)+(0.88,0.62)$) {$\theta_1=60^\circ$};
			}

			% Componentes F1
			\only<4->{
			\draw[component, blue] (O) -- (FxOne)
			node[midway, below=2pt, label, text=blue] {$F_{x1}$};

			\draw[component, blue] (FxOne) -- (Fone)
			node[midway, right=2pt, label, text=blue] {$F_{y1}$};
			}

			% Vector F2
			\only<5->{
			\draw[vector, green] (O) -- (Ftwo)
			node[pos=0.62, above left=2pt, label, text=green] {$\vec{F}_2$};

			\node[
				label,
				text=green,
				anchor=east
			] at ($(Ftwo)+(-0.18,0.08)$) {$F_2=70$};
			}

			% Ángulo theta2
			\only<6->{
			\draw[green, thick]
			($(O)+(1.45,0)$)
			arc[start angle=0,end angle=113,radius=1.45];

			\node[
				label,
				text=green,
				anchor=south
			] at ($(O)+(-0.10,1.62)$) {$\theta_2=113^\circ$};
			}

			% Componentes F2
			\only<7->{
			\draw[component, green] (O) -- (FxTwo)
			node[midway, below=2pt, label, text=green] {$F_{x2}$};

			\draw[component, green] (FxTwo) -- (Ftwo)
			node[midway, left=2pt, label, text=green] {$F_{y2}$};
			}

			% -------------------------------------------------------------
			% Panel derecho: tabla de componentes
			% -------------------------------------------------------------
			\only<8->{
			\node[
				panel,
				minimum width=10.15cm,
				minimum height=9.25cm,
				anchor=north west
			] at (2.35,6.85) {};

			\node[
				anchor=west,
				font=\bfseries\small,
				text=navy
			] at (2.75,6.35) {Tabla de componentes};

			\node[
			eq,
			anchor=north west,
			text width=9.15cm
			] at (2.80,5.75) {
			\[
			\begin{array}{c|c|c}
				\textbf{Vector} & F_x & F_y \\
				\hline
				\vec{F}_1       &
				F_1\cos\theta_1 &
				F_1\sin\theta_1             \\[1mm]
				\vec{F}_2       &
				F_2\cos\theta_2 &
				F_2\sin\theta_2
			\end{array}
		\]
		};
		}

		\only<9->{
				\node[
					eq,
					anchor=north west,
					text width=9.15cm
				] at (2.80,3.55) {
					Para $\vec{F}_1$:
					\[
						\begin{aligned}
							F_{x1} & = 50\cos(60^\circ)=25.00 \\[1mm]
							F_{y1} & = 50\sin(60^\circ)=43.30
						\end{aligned}
					\]
				};
			}

		\only<10->{
				\node[
					eq,
					anchor=north west,
					text width=9.15cm
				] at (2.80,1.30) {
					Para $\vec{F}_2$:
					\[
						\begin{aligned}
							F_{x2} & = 70\cos(113^\circ)=-27.35 \\[1mm]
							F_{y2} & = 70\sin(113^\circ)=64.44
						\end{aligned}
					\]
				};
			}

	\end{tikzpicture}

\end{frame}


% -------------------------------------------------------------
% Frame 2: sumatorias, resultante y DCL
% -------------------------------------------------------------
\begin{frame}{Sumatoria de vectores: resultante y ángulo}

	\centering

	\begin{tikzpicture}[
			x=1cm,
			y=1cm,
			>=Stealth,
			vector/.style={-{Stealth[length=3.2mm,width=2.4mm]}, very thick, line cap=round},
			component/.style={-{Stealth[length=2.8mm,width=2.1mm]}, very thick, dashed, line cap=round},
			axis/.style={-{Stealth[length=2.6mm,width=2.0mm]}, thick, muted},
			panel/.style={rounded corners=4pt, draw=muted!35, fill=paper, inner sep=5pt},
			label/.style={font=\scriptsize, inner sep=1pt},
			eq/.style={font=\footnotesize, align=left, text=ink},
			every node/.style={font=\small}
		]

		% =============================================================
		% Datos finales:
		% Sum Fx = -2.35
		% Sum Fy = 107.74
		% R = 107.76
		% theta_R = 91.25°
		% =============================================================

		% -------------------------------------------------------------
		% Panel izquierdo: DCL de la resultante
		% -------------------------------------------------------------
		\only<1->{
				\node[
					panel,
					minimum width=12.4cm,
					minimum height=9.25cm,
					anchor=north west
				] at (-13.15,6.85) {};

				\node[
					anchor=west,
					font=\bfseries\small,
					text=navy
				] at (-12.70,6.35) {DCL de la resultante};
			}

		\coordinate (Q) at (-7.15,1.10);

		% Ejes DCL
		\only<1->{
				\draw[axis] ($(Q)+(-3.0,0)$) -- ($(Q)+(3.6,0)$)
				node[right,label,text=ink] {$+x$};

				\draw[axis] ($(Q)+(0,-2.0)$) -- ($(Q)+(0,4.7)$)
				node[above,label,text=ink] {$+y$};

				\node[label,text=muted] at ($(Q)+(-0.18,-0.25)$) {$O$};
			}

		% Componente en x
		\only<2->{
				\draw[component, red] (Q) -- ++(-0.90,0)
				node[left,label,text=red] {$\sum F_x=-2.35$};
			}

		% Componente en y
		\only<3->{
				\draw[component, red] (Q) -- ++(0,3.85)
				node[above,label,text=red] {$\sum F_y=107.74$};
			}

		% Resultante
		\only<4->{
				\draw[vector, red] (Q) -- ++(-0.90,3.85)
				node[above left=2pt,label,text=red] {$\vec{R}$};

				\node[
					label,
					text=red,
					anchor=west
				] at ($(Q)+(-0.45,4.10)$) {$R=107.76$};
			}

		% Ángulo correcto medido desde +x
		\only<5->{
				\draw[red, thick]
				($(Q)+(1.05,0)$)
				arc[start angle=0,end angle=91.25,radius=1.05];

				\node[
					label,
					text=red,
					anchor=west
				] at ($(Q)+(0.72,1.05)$) {$\theta_R=91.25^\circ$};
			}

		% Nota conceptual
		\only<6->{
				\node[
					eq,
					anchor=north west,
					text width=9.3cm
				] at (-12.55,-1.85) {
					La resultante queda casi vertical, ligeramente hacia la izquierda, porque:
					\[
						\sum F_x<0,
						\qquad
						\sum F_y>0
					\]
				};
			}

		% -------------------------------------------------------------
		% Panel derecho superior: sumatorias
		% -------------------------------------------------------------
		\only<7->{
				\node[
					panel,
					minimum width=10.15cm,
					minimum height=4.85cm,
					anchor=north west
				] at (2.35,6.85) {};

				\node[
					anchor=west,
					font=\bfseries\small,
					text=navy
				] at (2.75,6.35) {Sumatorias};

				\node[
					eq,
					anchor=north west,
					text width=9.15cm
				] at (2.80,5.70) {
					\[
						\begin{aligned}
							\sum F_x & = F_{x1}+F_{x2}  \\
							         & = 25.00+(-27.35) \\
							         & = \boxed{-2.35}
						\end{aligned}
					\]
				};
			}

		\only<8->{
				\node[
					eq,
					anchor=north west,
					text width=9.15cm
				] at (2.80,3.40) {
					\[
						\begin{aligned}
							\sum F_y & = F_{y1}+F_{y2}  \\
							         & = 43.30+64.44    \\
							         & = \boxed{107.74}
						\end{aligned}
					\]
				};
			}

		% -------------------------------------------------------------
		% Panel derecho inferior: magnitud y ángulo
		% -------------------------------------------------------------
		\only<9->{
				\node[
					panel,
					minimum width=10.15cm,
					minimum height=5.95cm,
					anchor=north west
				] at (2.35,1.55) {};

				\node[
					anchor=west,
					font=\bfseries\small,
					text=navy
				] at (2.75,1.05) {Vector resultante};

				\node[
					eq,
					anchor=north west,
					text width=9.15cm
				] at (2.80,0.45) {
					\[
						\begin{aligned}
							R & = \sqrt{(\sum F_x)^2+(\sum F_y)^2} \\
							  & = \sqrt{(-2.35)^2+(107.74)^2}      \\
							  & = \boxed{107.76}
						\end{aligned}
					\]
				};
			}

		\only<10->{
				\node[
					eq,
					anchor=north west,
					text width=9.15cm
				] at (2.80,-2.10) {
					\[
						\begin{aligned}
							\theta_R
							 & = \operatorname{atan2}\left(\sum F_y,\sum F_x\right) \\
							 & = \operatorname{atan2}(107.74,-2.35)                 \\
							 & = \boxed{91.25^\circ}
						\end{aligned}
					\]
				};
			}

	\end{tikzpicture}

\end{frame}

\end{document}